\section{Spectral effectivity and the full Tate conjecture: the precise numerical-motive equivalence}

The SA06 intake, row D.4f and its successor-work instruction at
lines 501--503, asks whether effectivity detected by Frobenius is
equivalent to the Tate conjecture for all motives.  Its earlier proof
file is absent, but this question has the following precise answer in
the category of rational numerical motives.  This is a deduction from
classical results, not a new solution of the Tate conjecture.

Fix $k=\mathbb F_q$.  Let $\mathcal N(k)$ be the pseudo-abelian
category of pure numerical motives with rational coefficients and all
integer Tate twists.  Call an object \emph{effective} if it is a
direct summand of a finite direct sum of motives $h(X)$ with $X/k$
smooth projective, without a positive Tate twist.  The Frobenius
endomorphism $\pi$ is central.  On a simple object $S$, the algebra
$\mathbb Q[\pi_S]$ is a number field, so $\pi_S$ defines an algebraic
number up to conjugacy.  Call $M$ \emph{spectrally integral} if all
these numbers for its simple constituents are algebraic integers.
This definition does not presume that an $\ell$-adic realization
factors through numerical equivalence.

The structural inputs are Jannsen's semisimplicity theorem and
Frobenius centrality, in the exact formulations of Geisser (1998),
Theorem 2.1, p.~111, and Propositions 2.3--2.4, p.~112.
The decisive known result is Geisser, Theorem 1.1, p.~109,
equivalently Theorem 2.7, p.~114: the full Tate conjecture over $k$
is equivalent to the assertion that a simple numerical motive whose
Frobenius is $1$ is the unit.  We use this theorem as a cited input.
Here \emph{full Tate} means equality of the pole order at $s=r$
and the dimension of codimension-$r$ numerical cycle classes, for
all smooth projective $X/k$ and all $r$; it is not merely the
surjectivity assertion for one variety and one degree.

\subsection{BR-EFF01. The equivalence and its weight-zero test}
The following assertions are equivalent:
\begin{enumerate}
\item The full Tate conjecture holds over $k$.
\item Every spectrally integral object of $\mathcal N(k)$ is effective.
\item Every object $M$ of $\mathcal N(k)$ with $\pi_M=1_M$ is effective.
\item Every simple object $S$ with $\pi_S=1_S$ is effective.
\end{enumerate}
For every effective object, spectral integrality holds without
any Tate-conjecture assumption.  Thus item 2 can equivalently be
written ``effective if and only if spectrally integral.''

\paragraph{Complete deduction.}
For smooth projective $X$, its cohomological Frobenius characteristic
polynomials are monic integral polynomials, and their product annihilates Frobenius
in the category of homological motives.  They therefore annihilate
it numerically.  The minimal polynomial on a numerical summand
divides such a polynomial, so all its roots are algebraic integers.
This proves the unconditional necessity.

Assume item 4, and let $S$ be simple with $\pi_S=1$.  Since $S$
is effective, it is a summand of some finite sum of $h(X)$.
The algebraic K\"unneth projectors over $\mathbb F_q$ give
$h(X)=\bigoplus_{i\geq0}h^i(X)$; the precise cited input is
Katz--Messing (1974), Theorem 2(1), pp.~76--77.
For $i>0$, a polynomial annihilating $\pi_{h^i(X)}$ has all roots
of absolute value $q^{i/2}>1$.  In particular $\pi_{h^i(X)}-1$
is invertible: a B\'ezout identity between its annihilating
polynomial and $T-1$ constructs the inverse.  Centrality of
Frobenius now forces every map $S\to h^i(X)$ and $h^i(X)\to S$
to be zero.  Hence $S$ is a summand of a finite sum of $h^0(X)$.

These degree-zero objects are Artin motives: $h^0(X)$ is the
motive of the finite \'etale scheme of geometric connected
components.  For a geometrically connected component the
degree-zero projector is obtained from a degree-one rational
zero-cycle, obtained by dividing any closed point by its degree.
For a component with constant field $\mathbb F_{q^a}$ this gives
$h(\operatorname{Spec}\mathbb F_{q^a})$.  In its permutation
representation Frobenius is a cyclic permutation, and the
$1$-eigenspace is the constant line.  The rational projector is
$a^{-1}\sum_{j=0}^{a-1}\pi^j$.  Thus the part with Frobenius $1$
in any finite sum of such Artin objects is a sum of units.  The
simple object $S$ is consequently the unit.  Geisser's cited
theorem proves item 1.

Suppose now that item 1 holds.  The same cited theorem says that
simple numerical motives are determined by their Frobenius number.
Let $S$ be simple and spectrally integral, with Frobenius number
$\alpha$.  Its weight is an integer $m$, and all conjugates have
absolute value $q^{m/2}$.  Because $\alpha$ is a nonzero algebraic
integer, the norm shows $m\geq0$.

If $m=0$, Kronecker's theorem says that $\alpha$ is a root of
unity.  The rational Artin representation for its cyclotomic
polynomial is a summand of the regular representation of a finite
cyclic quotient of $\operatorname{Gal}(\overline k/k)$, hence is
an effective Artin motive.  Its simple constituent with Frobenius
$\alpha$ is isomorphic to $S$ by Geisser's characterization.

If $m>0$, choose an algebraic root $\beta$ with $\beta^m=\alpha$.
It is integral and each conjugate has absolute value $q^{1/2}$.
Honda--Tate supplies an abelian variety $B/k$ having $\beta$ among
its $H^1$ Frobenius eigenvalues.  The tensor power
$h^1(B)^{\otimes m}$ is an effective summand of $h(B^m)$, and its
realization has $\alpha$ among its eigenvalues.  Under full Tate,
numerical and homological equivalence agree and realization is
faithful.  Its decomposition into simple numerical summands
therefore contains a summand $S_\alpha$ with Frobenius number
$\alpha$.  Geisser's characterization gives $S\simeq S_\alpha$.
Thus $S$ is effective.  A finite direct sum of effective motives
is effective, proving item 2.  The implications $2\Rightarrow3
\Rightarrow4$ are immediate.  This proves all assertions.
\(\square\)

This proof never assumes a rational neutral fibre functor.
The quaternion obstruction in AQ-RATIONAL-FIBRE is compatible
with every step.  It also does not identify numerical and Chow
effectivity: a lift of a numerical projector to a Chow projector
requires additional information.

\subsection{BR-EFF02. The effective cone in the now established abelian sector}
Let $\mathcal A_q$ be the rational numerical category generated by
abelian, Artin, and Tate motives, with the same conventions as
CC-AG01.  Using the cited all-abelian Tate result in CC-TATE01--04,
an object of $\mathcal A_q$ is effective if and only if it is
spectrally integral.  In fact it then has an effective presentation
using only abelian varieties and finite \'etale schemes.  Over
$E=\overline{\mathbb Q}_\ell$, the equivalence of CC-AG01 identifies
this full subcategory with the finite $W(q)$-graded vector spaces
whose nonzero components are indexed by
\[
 W^+(q)=\{\alpha\in W(q):\alpha\text{ is an algebraic integer}\}.
 \tag{EFF.1}
\]
This is a tensor monoid of labels, not a group of labels.

\paragraph{Complete deduction.}
Necessity is the unconditional first paragraph of BR-EFF01.
For sufficiency, apply its construction to every simple constituent.
The required comparisons occur entirely inside $\mathcal A_q$.
Their validity follows from the all-abelian Tate theorem and its
full-faithfulness consequence AQ-ABELIAN-RECOGNITION: for simple
objects in this sector, an intersecting Frobenius spectrum gives a
nonzero morphism after scalar extension, hence a nonzero rational
morphism, and such a morphism is an isomorphism by semisimplicity.
In weight zero use the finite \'etale construction; in positive
weight use $h^1(B)^{\otimes m}$ as above.  These constructions
stay effective before any coefficient extension.

After extension to $E$, all rank-one $W(q)$-labels split by
CC-AG01.  Each label in $W^+(q)$ is effective by the same
construction, and each effective label is integral by necessity.
Finite sums and summands give (EFF.1).  Multiplication preserves
integrality, so $W^+(q)$ is a monoid.  For example $q\in W^+(q)$
whereas $q^{-1}\notin W^+(q)$, so it is not a group.
\(\square\)

\subsection{BR-EFF03. The slope test and the essential Artin correction}
For a Weil $q$-number $\alpha$ of weight $m$, spectral integrality
is equivalent to
\[
 \frac{v(\alpha)}{v(q)}\geq0
 \quad\text{for every finite place }v\mid p
 \text{ of a number field containing }\alpha.
 \tag{EFF.2}
\]
Consequently (EFF.2) is an exact effectivity test in the abelian
sector of BR-EFF02.  For all numerical motives the corresponding
universal assertion is equivalent to full Tate by BR-EFF01.

\paragraph{Complete proof.}
The definition of a Weil number gives $q^N\alpha$ integral for
some $N$.  At places away from $p$, this implies $v(\alpha)\geq0$.
Choose a normal field containing $\alpha$ and its conjugates.
For a complex-conjugation automorphism $c$ one has
$\alpha c(\alpha)=q^m$.  The same nonnegativity holds for
$c(\alpha)$ away from $p$, and their valuations sum to zero
there.  Both valuations are therefore zero.  Hence the only
possible denominators of $\alpha$ occur over $p$, and (EFF.2)
is precisely their absence.  This proves the equivalence.
\(\square\)

There is a small but necessary source correction in using an
effective realization theorem.  Milne--Ramachandran (2009),
Proposition 1.7 in the author final version of May 2, 2008,
p.~4, is printed for every effective semisimple Tate structure
with an abelian variety as its sole source.  Its unrestricted
weight-zero wording has the following counterexample.
The one-dimensional space with Frobenius $-1$ is effective,
semisimple, and of weight zero.  For any abelian variety $B/k$,
$H^0(B,\mathbb Q_\ell)$ has Frobenius $1$, whereas every
positive-degree eigenvalue has absolute value greater than one.
Thus $-1$ never occurs in $H^*(B,\mathbb Q_\ell)$.
The space does occur in
$H^0(\operatorname{Spec}\mathbb F_{q^2},\mathbb Q_\ell)$, as
the anti-invariant summand.  This proves the counterexample and
the required repair in that case.  The general weight-zero repair
is the finite cyclic Artin construction in BR-EFF01; the positive
weight construction there is complete.  The effective targets of
CC-XC01 already include these Artin summands.  No refutation of
the generalized Tate conjecture is inferred from this wording
correction.

\paragraph{Exact sources.}
\begin{itemize}
\item Geisser, T. (1998). Tate's conjecture, algebraic cycles and
rational $K$-theory in characteristic $p$. \emph{$K$-Theory,
13}(2), 109--122.
\url{https://www2.rikkyo.ac.jp/web/geisser/Tate_KtheoryJournal.pdf}.
Theorem 1.1, p.~109; Theorem 2.1 and Corollary 2.2, p.~111;
Propositions 2.3--2.4, p.~112; Theorem 2.7, p.~114.
\item Milne, J. S., \& Ramachandran, N. (2009). Motivic complexes
over finite fields and the ring of correspondences at the generic
point. \emph{Pure and Applied Mathematics Quarterly, 5}(4),
1219--1252. Operative author final version, May 2, 2008:
\url{https://www.jmilne.org/math/articles/2009a.pdf}.
Proposition 1.7, p.~4; Remark 1.12, p.~6; Lemma 2.1 and
Proposition 2.2, p.~8. The weight-zero correction above is explicit.
\item Milne, J. S. (2007, September 19). \emph{The Tate conjecture
over finite fields (AIM talk)}.
\url{https://www.jmilne.org/math/xnotes/TateAimv1.pdf}.
Section 1, pp.~1--4, particularly footnote 4, p.~4, on motivic
effectivity and generalized Tate. This contextual statement is
not substituted for the precise numerical-motive theorem proved here.
\end{itemize}
